\clearpage
\SourceLeaf{172}\centerline{[ZUR THEORIE DER PARALLELLINIEN.]}

\bigskip
\hrule
\bigskip

\centerline{[1.]}

\smallskip

\centerline{\textsc{Gauss} an \textsc{Olbers}. \ Braunschweig, 30. Juli 1806.}

\smallskip

.\ .\ .\ .\ .\ Herr \textsc{von Zach} schreibt mir noch, dass Sie sich angeboten hätten, \textsc{Legendres} Werk über die Kometenbahnen zu rezensieren. Mit Vergnügen werde ich Ihnen also das mir von Herrn \textsc{von Zach} zugeschickte Exemplar nach Bremen senden, doch erlauben Sie wohl, dass ich es noch einige Wochen behalte. Beim ersten Durchblättern scheint es mir sehr viel Schönes zu enthalten. Vieles von dem, was an meiner Methode besonders in ihrer ersten Gestalt eigentümlich war, finde ich auch in diesem Buch wieder. Es scheint mein Schicksal zu sein, in fast allen meinen theoretischen Arbeiten mit \textsc{Legendre} zu konkurrieren. So in der höheren Arithmetik, in den Untersuchungen über transzendente Funktionen, die mit der Rektifikation der Ellipse zusammenhängen, bei den ersten Gründen der Geometrie und nun wieder hier. So wird z.\,B. auch das von mir seit 1794 gebrauchte Prinzip, dass man, um mehrere Größen, die man nicht alle genau darstellen kann, möglichst gut darzustellen, die Summe der Quadrate zu einem Minimum machen muss, in \textsc{Legendres} Werk verwendet und recht geschickt ausgeführt.\ .\ .\ .\ .\ .\ .

\bigskip
\hrule
\bigskip

\centerline{[2.]}

\smallskip

\centerline{[\emph{Aus} \textsc{Schumachers} \emph{Tagebuch: »Gaussiana«. Göttingen, November 1808.}]}

\smallskip

\{\textsc{Gauss} hat die Theorie der Parallellinien darauf zurückgeführt, dass es, wenn die angenommene Theorie nicht wahr wäre, eine konstante Linie geben müsste, deren Länge von vornherein gegeben wäre, was absurd ist. Doch hält er selbst diese Arbeit noch nicht für hinreichend.\}

\bigskip
\hrule
